\section{The finite complement and completeness}
\label{sec:finite}

The preceding argument covers every even normalized parameter
\(A\le-13\) and every odd normalized parameter \(A\le-17\).
The remaining cases are finite for a proved reason. We give their
complete certificates, including an exact procedure for checking the
point sets rather than only the number of points.

\begin{lemma}[Finite partial permutations]
\label{lem:pruning}
Let \(F:X\to X\) be injective, and suppose every periodic point of
\(F\) belongs to a specified finite set \(V_0\subset X\).
Define
\begin{equation}
 V_{j+1}=\{p\in V_j:F(p)\in V_j\}.
 \label{eq:pruning}
\end{equation}
This decreasing sequence stabilizes at a set \(V_*\), which is exactly
the periodic set of \(F\) in \(X\).
\end{lemma}
\begin{proof}
Every periodic orbit contained in \(V_0\) is contained in each \(V_j\)
by induction. A decreasing sequence of finite sets stabilizes, and the
restriction \(F:V_*\to V_*\) is injective. It is therefore a
permutation of \(V_*\), so every point of \(V_*\) is periodic.
\end{proof}

\subsection{The thirteen even parameters}

For \(c=0,\ldots,12\), so that \(A=-c\), use the integer quantities
\begin{equation}
 B_c=1+\lfloor\sqrt{1+c}\rfloor,\qquad
 S_c=\{u\in\Z:|u|\le B_c,\ |u^2-c|\le2B_c\}.
 \label{eq:even_finite_alphabet}
\end{equation}
Lemma~\ref{lem:real_bounds} shows that every periodic coordinate belongs
to \(S_c\). Start with \(V_0=S_c\times S_c\), and apply
Lemma~\ref{lem:pruning} to
\begin{equation}
 F_c(x,y)=(y,y^2-c-x).
 \label{eq:even_finite_map}
\end{equation}
The bound and alphabet use only an integer square root and integer
comparisons. Table~\ref{tab:even_certificate} gives every successive
cardinality up to stabilization and every stable coordinate word.

\subsection{The seventeen odd parameters}

For \(A=0,-1,\ldots,-16\), use doubled centered coordinates
\(z_i=2x_i+1\), which are odd integers. The bound
in~\eqref{eq:real_bounds} gives
\[
 |z_i|\le2+\sqrt{1-4A}.
\]
Define
\begin{equation}
 \begin{split}
 B_A&=2+\lfloor\sqrt{1-4A}\rfloor,\\
 S_A&=\{z\in\Z:z\text{ odd},\ |z|\le B_A,
                         \ |z^2+4A+3|\le4B_A\}.
 \end{split}
 \label{eq:odd_finite_alphabet}
\end{equation}
Indeed, the doubled recurrence is
\begin{equation}
 2(z_{i-1}+z_{i+1})=z_i^2+4A+3,
 \label{eq:doubled_recurrence}
\end{equation}
which proves the second alphabet constraint. On pairs of odd integers
the map is
\begin{equation}
 G_A(z,w)=\left(w,\frac{w^2+4A+3}{2}-z\right).
 \label{eq:odd_finite_map}
\end{equation}
For odd \(w\), the integer \(w^2+4A+3\) is divisible by four.
The quotient by two is even, so the second coordinate of \(G_A\) is
odd. Solving for \(z\) gives an inverse on the odd lattice, in
particular injectivity. Take \(V_0=S_A\times S_A\) and use
\eqref{eq:pruning}. Table~\ref{tab:odd_certificate} records the entire
calculation; its words are expressed in the original normalized
integer coordinates \(x=(z-1)/2\).

\subsection{What the certificates assert}

For each row of the appendix tables, let \(W\) be the displayed set of
coordinate words. Define the corresponding pair set by
\begin{equation}
 E(W)=\bigcup_{(v_0,\ldots,v_{d-1})\in W}
       \{(v_i,v_{i+1}):0\le i<d\},
 \label{eq:word_expansion}
\end{equation}
where \(v_d=v_0\). Each even row asserts \(V_*=E(W)\). Each odd row
asserts
\[
 V_*=\{(2x+1,2y+1):(x,y)\in E(W)\}.
\]
Thus the claimed final set is specified explicitly, not inferred from
its cardinality.

For completeness, the pruning calculation in every row is the
following exact finite procedure. Here \(S\) and \(F\) mean the
appropriate alphabet and map just defined.
\begin{quote}
\begin{verbatim}
V = {(x, y) : x in S and y in S}
sizes = [cardinality(V)]
repeat:
    U = {p in V : F(p) in V}
    if U = V:
        return V, sizes
    V = U
    append cardinality(V) to sizes
\end{verbatim}
\end{quote}
All sets consist of integer pairs. The division in \(G_A\) is exact,
as proved above. The last cardinality listed in a table row is unchanged
by one further step; a final zero denotes the empty stable set.
Expanding the words by~\eqref{eq:word_expansion} gives the full-set
comparison for each returned \(V\). The starting sets have at most
81 elements, so the certificate is finite without an orbit-period
cutoff.

The complete point sets were also verified by a separate construction:
in the original integer coordinates, form the entire bounding square
or rectangle, retain exactly those one-step edges that stay within it,
and compute Boolean transitive closure. A vertex is on a nonempty
return path precisely when it is periodic. This second construction
does not use the filtered alphabet or the pruning operation, and the
comparison includes all vertices and cycle words, not just totals.
The computational part is limited to the stated thirteen and seventeen
parameters; the all-parameter argument is
Section~\ref{sec:six_symbols} and Lemma~\ref{lem:symbol_classification}.

\begin{proof}[Completion of the classification in
Theorem~\ref{thm:classification}]
Lemma~\ref{lem:integrality} first reduces rational periodic points to
integer points, and Lemma~\ref{lem:translation} reduces the coefficients
to the two normal forms. Equations~\eqref{eq:sum_even} and
\eqref{eq:sum_odd} settle the upper parameter ranges. The two
six-symbol lemmas and Lemma~\ref{lem:symbol_classification} settle
every \(A\le-13\) in the even branch and every \(A\le-17\) in the
odd branch. The exact certificates in
Tables~\ref{tab:even_certificate} and~\ref{tab:odd_certificate}
settle the complementary parameters. Every word there occurs in
Tables~\ref{tab:even} or~\ref{tab:odd}, and
Lemma~\ref{lem:existence} proves existence and exact periods for all
entries of those tables. This proves the complete orbit lists and the
period restriction. The count and equality locus are proved next.
\end{proof}
